\begin{abstract}
The fastest parallel implementations of breadth-first search (BFS) for shared memory are direction-optimizing and blocking: every level or round ends at a barrier or a join, so a thread that is preempted or stopped by a page fault delays all the others. A BFS is wait-free if every call returns after a finite number of the caller's own steps, which is easy if each thread searches the whole graph alone. The challenge is to share a search among threads without waiting, while repeating little work and competing in speed with the best blocking implementations.

We present such a method. The vertices that a thread discovers are its share of the next level; it expands its own share, then helps with the rest, and every write that lets other threads skip work comes after that work, so a thread that stops leaves nothing that the others cannot redo. We prove the method correct and wait-free under the memory model of x86, given that allocation takes a bounded number of steps, bound the work of a top-down level by that level and the number of threads, and give a variant whose memory depends on the largest levels, plus a constant per level.

On five standard graphs and 128 cores, our method examines about as many neighbors as the sequential baseline at 16 threads, and at most 30\% more at 112; the best times of GAPBS, GBBS and PASGAL are 1.58, 2.77 and 2.73 times its own in the geometric mean. With one of 112 threads stalled for 10~ms at every level, or 1~ms on the road network, its median time grows by at most 7\%, and theirs up to 19.5-fold. Since every thread takes part in every level, it is slow on long paths and other graphs with many tiny levels.
\end{abstract}
